\chapter{Tempered Distributions and the Generalized Transform}
\label{ch:tespered}

\begin{quote}
\itshape
The transform cannot be transplanted to $\D'$ because it does not preserve $\D$. It is preserved by
a smaller but more regular space, the Schwartz space, and on the dual of that space the transform is
defined by moving it to the probe. Inversion is proved independently by Gaussian regularization, so
nothing circular enters. The three results of the preceding chapters then become theorems.
\end{quote}

\section{Why the naive definition fails}

The natural attempt is
\begin{equation}
\ip{\F T}{\varphi}:=\ip{T}{\F\varphi},
\label{eq:naive}
\end{equation}
which requires $\F\varphi$ to be a test function whenever $\varphi$ is. It is not.

\begin{proposition}\label{prop:FnotD}
If $\varphi\in C_{c}^{\infty}$ and $\varphi\not\equiv0$, then $\F\varphi\notin C_{c}^{\infty}$.
\end{proposition}

\begin{proof}
For complex $\zeta$ the integral $\ph(\zeta)=\int\varphi(t)e^{-i\zeta t}\dd t$ converges absolutely
because $\varphi$ is compactly supported, and it defines an entire function of $\zeta$. An entire
function vanishing on a nonempty open set vanishes identically, whence $\varphi\equiv0$ by inversion.
So $\ph$ cannot have compact support unless $\varphi\equiv0$.
\end{proof}

\begin{example}
The rectangular window $w_{T}=\mathbf{1}_{[-T,T]}$ has transform $2\sin(T\omega)/\omega$, which
oscillates along the whole real line. Smoothing the corners does not shorten its tail; by
Proposition~\ref{prop:FnotD} no nonzero smooth compactly supported $\varphi$ has a compactly supported
transform.
\end{example}

The conclusion is structural: a space of probes on which $\F$ acts must be invariant under the
transform, and $C_{c}^{\infty}$ is not.

\section{The Schwartz space}

\begin{definition}\label{def:S}
$\Sw(\R)$ is the set of $\varphi\in C^{\infty}(\R)$ such that for all $m,n\in\mathbb{N}$,
\begin{equation}
p_{m,n}(\varphi):=\sup_{x\in\R}(1+|x|)^{m}\big|\varphi^{(n)}(x)\big|<\infty .
\label{eq:seminorm}
\end{equation}
The $p_{m,n}$ are seminorms, and $\varphi_{j}\to\varphi$ in $\Sw$ means
$p_{m,n}(\varphi_{j}-\varphi)\to0$ for all $m,n$.
\end{definition}

Equivalently: every derivative decays faster than any negative power of $|x|$.

\begin{example}\label{ex:S}
$e^{-x^{2}}\in\Sw$; every element of $C_{c}^{\infty}$ belongs to $\Sw$; but $1$ and $x^{n}$ do not
(they do not decay), $e^{x^{2}}$ does not (it grows), and $1/(1+x^{2})$ does not, although it is
integrable: it decays only like $x^{-2}$, and its derivatives decay more slowly still. Integrability
and rapid decay are different requirements.
\end{example}

\section{The transform acts on the Schwartz space}

Two identities carry the whole proof.

\begin{lemma}\label{lem:two-identities}
For $\varphi\in\Sw$ and $n,m\in\mathbb{N}$,
\begin{equation}
\partial_{\omega}^{\,n}\ph(\omega)=\F\big[(-it)^{n}\varphi\big](\omega),
\qquad
\omega^{m}\,\ph(\omega)=i^{-m}\,\F\big[\varphi^{(m)}\big](\omega).
\label{eq:two-identities}
\end{equation}
\end{lemma}

\begin{proof}
The first is Lemma~\ref{lem:diffint} iterated. For the second, Theorem~\ref{thm:derivt} gives
$\F[\varphi^{(m)}]=(i\omega)^{m}\ph$, and multiplying by $i^{-m}$ yields \eqref{eq:two-identities}.
\end{proof}

\begin{theorem}\label{prop:F-S}
$\F$ maps $\Sw$ continuously into $\Sw$; that is, $\ph\in\Sw$ for $\varphi\in\Sw$, and for all
$m,n$ there are $C$ and finitely many indices $(m_{k},n_{k})$ with
$p_{m,n}(\ph)\le C\sum p_{m_{k},n_{k}}(\varphi)$.
\end{theorem}

\begin{proof}
Combining the two identities,
\begin{equation}
\omega^{m}\partial_{\omega}^{\,n}\ph(\omega)
=i^{-m}\int_{\R}\partial_{t}^{m}\big(t^{n}\varphi(t)\big)e^{-i\omega t}\dd t ,
\end{equation}
hence
\begin{equation}
\big|\omega^{m}\partial_{\omega}^{\,n}\ph(\omega)\big|
\le\big\|\partial_{t}^{m}(t^{n}\varphi)\big\|_{1} .
\label{eq:Fbound}
\end{equation}
The right side does not depend on $\omega$. Since $\varphi\in\Sw$, the function $t^{n}\varphi$ lies in
$\Sw$, so each of its derivatives is integrable, and the right side of \eqref{eq:Fbound} is finite.
Taking the supremum over $\omega$ and using $(1+|\omega|)^{m}\le C_{m}\sum_{k\le m+1}|\omega|^{k}$
yields the stated bound.
\end{proof}

The proof shows why $\Sw$ is the right space. The two operations that appear in the transform---taking
a derivative and multiplying by the variable---exchange places under $\F$, and $\Sw$ is closed under
both.

\section{Inversion}

The inversion formula cannot be obtained from
$\int e^{i\omega(t-s)}\dd\omega=2\pi\delta(t-s)$ without circularity: that identity is itself a
distributional statement about $\delta$. The following proof uses only the Gaussian integral, Fubini,
and the approximate identities of Chapter~\ref{ch:limits}.

\begin{lemma}[Gaussian integral]\label{lem:gauss}
For $\eta>0$ and $x\in\R$,
\begin{equation}
\int_{\R}e^{-\eta\omega^{2}/2}e^{i\omega x}\dd\omega=\sqrt{\frac{2\pi}{\eta}}\;e^{-x^{2}/(2\eta)} .
\label{eq:gaussint}
\end{equation}
\end{lemma}

\begin{proof}
Complete the square in the exponent and use $\int e^{-u^{2}/2}\dd u=\sqrt{2\pi}$.
\end{proof}

\begin{theorem}[Fourier inversion]\label{thm:inversion}
Let $\varphi\in\Sw$. Then
\begin{equation}
\varphi(t)=\frac{1}{2\pi}\int_{\R}\ph(\omega)e^{i\omega t}\dd\omega
\qquad\text{for every }t\in\R .
\label{eq:inversion}
\end{equation}
\end{theorem}

\begin{proof}
For $\eta>0$ define
\begin{equation}
I_{\eta}(t)=\frac{1}{2\pi}\int_{\R}\ph(\omega)e^{i\omega t}e^{-\eta\omega^{2}/2}\dd\omega .
\end{equation}
Insert $\ph(\omega)=\int\varphi(s)e^{-i\omega s}\dd s$. The double integral is absolutely convergent,
because the modulus of the integrand is $|\varphi(s)|e^{-\eta\omega^{2}/2}$ and $\varphi\in\Lone$,
while the Gaussian is integrable in $\omega$. Fubini therefore applies and
\begin{align}
I_{\eta}(t)
&=\frac{1}{2\pi}\int_{\R}\varphi(s)\left[\int_{\R}e^{-\eta\omega^{2}/2}e^{i\omega(t-s)}\dd\omega
\right]\dd s \notag\\
&=\frac{1}{2\pi}\int_{\R}\varphi(s)\sqrt{\frac{2\pi}{\eta}}\,e^{-(t-s)^{2}/(2\eta)}\dd s
=\big(\rho_{\eta}*\varphi\big)(t),
\label{eq:conv}
\end{align}
where $\rho_{\eta}(u)=(2\pi\eta)^{-1/2}e^{-u^{2}/(2\eta)}$ is the Gaussian approximate identity.

Two limits are now taken. On the left, since $\ph\in\Sw\subset\Lone$ and
$|e^{-\eta\omega^{2}/2}|\le1$, dominated convergence with the $\eta$-independent bound $|\ph|$ gives
\begin{equation}
I_{\eta}(t)\longrightarrow\frac{1}{2\pi}\int_{\R}\ph(\omega)e^{i\omega t}\dd\omega .
\end{equation}
On the right, Theorem~\ref{thm:conv} gives $(\rho_{\eta}*\varphi)(t)\to\varphi(t)$. Equating the two
limits yields \eqref{eq:inversion}.
\end{proof}

\begin{theorem}[two transforms give a reflection]\label{thm:F2}
For $\varphi\in\Sw$,
\begin{equation}
\F^{2}\varphi(t)=2\pi\,\varphi(-t).
\label{eq:F2}
\end{equation}
\end{theorem}

\begin{proof}
By definition $\F^{2}\varphi(t)=\int_{\R}\ph(\omega)e^{-i\omega t}\dd\omega$, and
\eqref{eq:inversion} with $t$ replaced by $-t$ identifies this with $2\pi\varphi(-t)$.
\end{proof}

\begin{remark}
The order of the argument matters. Inversion is proved first, without reference to $\delta$; only
afterwards may one use $\F^{2}=2\pi\,(\text{reflection})$ to obtain statements about $\delta$ cheaply.
Using the distributional identity $\int e^{i\omega t}\dd\omega=2\pi\delta(t)$ to prove inversion would
reverse the implication.
\end{remark}

\section{Tempered distributions and the dual definition}

\begin{definition}\label{def:Sprime}
A \emph{tempered distribution} is a continuous linear functional on $\Sw(\R)$. The space is denoted
$\Sw'(\R)$. Continuity is equivalent to the existence of $C<\infty$ and finitely many indices such
that
\begin{equation}
|\ip{T}{\varphi}|\le C\sum_{m\le M,\ n\le N}p_{m,n}(\varphi)
\qquad\text{for all }\varphi\in\Sw .
\label{eq:Sprime-est}
\end{equation}
\end{definition}

\begin{example}\label{ex:tempered}
Bounded functions are tempered, since $|\ip{f}{\varphi}|\le\|f\|_{\infty}\|\varphi\|_{1}$ and
$\|\varphi\|_{1}$ is controlled by $p_{2,0}$. Functions of polynomial growth are tempered, with the
index $N$ equal to the growth order plus $2$. $\delta^{(n)}$ is tempered, since
$|\ip{\delta^{(n)}}{\varphi}|=|\varphi^{(n)}(0)|\le p_{0,n}(\varphi)$. The principal value is
tempered, the estimate near the origin using one derivative of $\varphi$ and the estimate at
infinity using rapid decay.
\end{example}

\begin{proposition}[not every distribution is tempered]\label{prop:ex2}
The function $e^{x^{2}}$ is locally integrable, hence defines a distribution, but it is not a
tempered distribution.
\end{proposition}

\begin{proof}
Take $0\not\equiv\psi\in C_{c}^{\infty}$ nonnegative and set $\psi_{R}(x)=\psi(x-R)$. The seminorms
$p_{m,n}(\psi_{R})$ grow at most polynomially in $R$, because translation only multiplies the weight
$(1+|x|)^{m}$ by a factor $O(R^{m})$. On the other hand
\begin{equation}
\ip{e^{x^{2}}}{\psi_{R}}=\int e^{x^{2}}\psi(x-R)\dd x\ge e^{(R-r)^{2}}\int\psi ,
\end{equation}
where $r$ is the radius of $\supp\psi$, which grows faster than any polynomial. No estimate of the
form \eqref{eq:Sprime-est} can hold.
\end{proof}

\begin{definition}[generalized Fourier transform]\label{def:FTdist}
For $T\in\Sw'$ define $\F T\in\Sw'$ by
\begin{equation}
\ip{\F T}{\varphi}:=\ip{T}{\F\varphi},
\qquad \varphi\in\Sw .
\label{eq:dualdef}
\end{equation}
\end{definition}

\begin{proposition}[consistency with the classical transform]\label{prop:dual-consistency}
If $f,\varphi\in\Sw$, then $\ip{\F f}{\varphi}=\ip{f}{\F\varphi}$; consequently \eqref{eq:dualdef}
extends the classical transform.
\end{proposition}

\begin{proof}
By Fubini, justified by absolute convergence,
\begin{equation}
\ip{\F f}{\varphi}
=\iint f(t)e^{-i\omega t}\varphi(\omega)\dd t\dd\omega
=\int_{\R}f(t)\Big[\int_{\R}\varphi(\omega)e^{-i\omega t}\dd\omega\Big]\dd t
=\int_{\R}f(t)\ph(t)\dd t=\ip{f}{\F\varphi} . \qedhere
\end{equation}
\end{proof}

\begin{remark}[three details that must be checked]
First, \eqref{eq:dualdef} is well posed because $\varphi\in\Sw$ implies
$\F\varphi\in\Sw$ (Theorem~\ref{prop:F-S}). Second, the right side defines a continuous linear
functional of $\varphi$, hence an element of $\Sw'$. Third, there is no factor $1/(2\pi)$ and no
conjugate on the right; a formula with either present belongs to a different convention.
\end{remark}

\begin{theorem}[inverse transform]\label{thm:inverse}
$\F:\Sw'\to\Sw'$ is a linear isomorphism, with inverse
\begin{equation}
\F^{-1}T=\frac{1}{2\pi}\,\F\big(T(-\cdot)\big).
\end{equation}
\end{theorem}

\begin{proof}
By Theorem~\ref{thm:F2} in dual form, $\F^{2}T=2\pi\,T(-\cdot)$; solving gives the stated inverse.
\end{proof}

\section{Continuity, and the gaps it closes}

\begin{theorem}[the transform is continuous]\label{thm:F-cont}
If $T_{j}\to T$ in $\Sw'$, then $\F T_{j}\to\F T$ in $\Sw'$.
\end{theorem}

\begin{proof}
For fixed $\varphi\in\Sw$,
$\ip{\F T_{j}}{\varphi}=\ip{T_{j}}{\F\varphi}\to\ip{T}{\F\varphi}=\ip{\F T}{\varphi}$.
\end{proof}

This is the step missing from Chapters~\ref{ch:constant} and \ref{ch:step}, where a time-domain limit
and a frequency-domain limit were computed separately without a bridge between them.

\begin{example}[completing the constant]
For every $\varphi\in\Sw$, $|e^{-\ep|t|}\varphi(t)|\le|\varphi(t)|\in\Lone$ and
$e^{-\ep|t|}\varphi(t)\to\varphi(t)$ pointwise, so $e^{-\ep|\cdot|}\to1$ in $\Sw'$ by dominated
convergence. By Theorems~\ref{thm:F-cont} and \ref{prop:F-S},
\begin{equation}
\F\big[e^{-\ep|\cdot|}\big]\longrightarrow\F[1]
\qquad\text{in }\Sw',
\end{equation}
while Theorem~\ref{thm:exptransform} and Theorem~\ref{thm:sampling} give
$\F[e^{-\ep|\cdot|}]=G_{\ep}\to2\pi\delta$. Uniqueness of limits in $\Sw'$ yields
$\F[1]=2\pi\delta$. Compare \eqref{eq:F1}: the same statement, now proved rather than posited.
\end{example}

\section{The three results}

\begin{theorem}\label{thm:main}
In $\Sw'$:
\begin{align}
\F[\delta]&=1, \qquad \F[1]=2\pi\delta, \label{eq:Fdelta}\\
\F[t^{\,n}]&=2\pi i^{\,n}\delta^{(n)},\qquad n=0,1,2,\dots, \label{eq:Ftn}\\
\F[H]&=\pi\delta+\PV\frac{1}{i\omega}. \label{eq:FHgen}
\end{align}
\end{theorem}

\begin{proof}
\emph{(i) $\delta$ and $1$.} For $\varphi\in\Sw$,
\begin{equation}
\ip{\F\delta}{\varphi}=\ip{\delta}{\ph}=\ph(0)=\int_{\R}\varphi(t)\dd t=\ip{1}{\varphi},
\end{equation}
so $\F\delta=1$. For the second identity,
\begin{equation}
\ip{\F1}{\varphi}=\ip{1}{\ph}=\int_{\R}\ph(t)\dd t=2\pi\varphi(0)=\ip{2\pi\delta}{\varphi},
\end{equation}
where the middle equality is \eqref{eq:inversion} at $t=0$. This is the only place where inversion is
used, and it is the reason inversion had to be proved first.

\emph{(ii) The two differentiation rules.} First, from $\F[\varphi']=i\omega\ph$ evaluated at
$\omega=t$ one gets $(\varphi')^{\wedge}(t)=it\,\ph(t)$, that is
$t\,\ph(t)=-i(\varphi')^{\wedge}(t)$. Hence
\begin{equation}
\ip{\F(tT)}{\varphi}=\ip{T}{t\ph}
=-i\ip{T}{(\varphi')^{\wedge}}
=-i\ip{\widehat T}{\varphi'}
=i\ip{\widehat T'}{\varphi},
\end{equation}
so
\begin{equation}
\F(tT)=i\,\partial_{\omega}\widehat T .
\label{eq:multt-dist}
\end{equation}
Second, from $\ph'(\omega)=-i(t\varphi)^{\wedge}(\omega)$,
\begin{equation}
\ip{\F(T')}{\varphi}=\ip{T'}{\ph}=-\ip{T}{\ph'}
=i\ip{\widehat T}{t\varphi}
=i\ip{t\widehat T}{\varphi},
\end{equation}
so
\begin{equation}
\F(T')=i\omega\,\widehat T .
\label{eq:deriv-dist}
\end{equation}
Iterating \eqref{eq:multt-dist} on $\F[1]=2\pi\delta$ yields \eqref{eq:Ftn}.

\emph{(iii) The step.} $H_{\ep}=H\,e^{-\ep t}\to H$ in $\Sw'$ because
$|H_{\ep}\varphi|\le|\varphi|\in\Lone$; and by Theorem~\ref{thm:Heps} together with
Theorems~\ref{thm:realpart} and \ref{thm:imagpart},
$\F[H_{\ep}]\to\pi\delta-i\PV\frac{1}{\omega}$. Theorem~\ref{thm:F-cont} identifies the limit of the
transforms with $\F[H]$, and uniqueness gives \eqref{eq:FHgen}.
\end{proof}

\section{Other regularizations}

The theory must be independent of the chosen regularizing family. Two checks.

\subsection*{Gaussian}

For $g_{\sigma}(t)=e^{-\sigma t^{2}/2}$ one has, by the differential-equation method or
Lemma~\ref{lem:gauss},
\begin{equation}
\widehat{g_{\sigma}}(\omega)=\sqrt{\frac{2\pi}{\sigma}}\,e^{-\omega^{2}/(2\sigma)} .
\end{equation}
Its total mass is
$\sqrt{2\pi/\sigma}\cdot\sqrt{2\pi\sigma}=2\pi$, and it has width of order $\sqrt{\sigma}$ and height
of order $\sigma^{-1/2}$; hence $\widehat{g_{\sigma}}\to2\pi\delta$. The Gaussian route is entirely
independent of the exponential route of Chapter~\ref{ch:constant} and gives the same limit.

\subsection*{Rectangular windows, and sinc}

For $w_{T}=\mathbf{1}_{[-T,T]}$ one has $\widehat{w_{T}}(\omega)=2\sin(T\omega)/\omega$. Since
$w_{T}\to1$ in $\Sw'$ (dominated convergence with bound $|\varphi|$), Theorem~\ref{thm:F-cont} gives
\begin{equation}
\widehat{w_{T}}\longrightarrow\F[1]=2\pi\delta
\qquad\text{in }\Sw'.
\end{equation}

\begin{remark}[why Chapter~\ref{ch:limits} does not apply]
The sinc kernels are not nonnegative, they change sign infinitely often, and
$\int_{\R}\big|2\sin(T\omega)/\omega\big|\dd\omega$ grows like $\log T$. Neither the hypotheses of
Theorem~\ref{thm:sampling-general} nor those of Proposition~\ref{prop:general-kernels} hold. Their
convergence comes instead from the continuity of $\F$ applied to the convergent family $w_{T}\to1$.
There is no contradiction with the failure of pointwise convergence noted in
Proposition~\ref{prop:nolimit}: that failure concerns the pointwise behaviour of the truncated
integrals, this concerns distributional convergence of the transforms.
\end{remark}

\begin{remark}[what the extension does and does not do]
It does not assign values to divergent integrals. It does the following: it reinterprets a signal as
a functional, defines $\F$ on functionals by duality, proves that the definition is continuous and
invertible, and then observes that the three identities hold simultaneously. The extension is
conservative---it agrees with the classical transform wherever the latter exists---and it is
canonical, because inversion together with the pairing leaves no freedom in the definition.
\end{remark}

\section{Plancherel, and what the delta does not have}

The generalized transform extends the classical one. It does not extend Plancherel's identity to
$\delta$, and it is worth being precise about why.

\begin{theorem}[Plancherel]\label{thm:plancherel}
For $f,g\in\Sw$,
\begin{equation}
\int_{\R}f(t)\overline{g(t)}\dd t
=\frac{1}{2\pi}\int_{\R}\fh(\omega)\overline{\gh(\omega)}\dd\omega .
\label{eq:plancherel-pair}
\end{equation}
In particular $\|f\|_{2}^{2}=(2\pi)^{-1}\|\fh\|_{2}^{2}$, and $\F$ extends uniquely to a bounded
operator $L^{2}(\R)\to L^{2}(\R)$ with $\|\F f\|_{2}=\sqrt{2\pi}\,\|f\|_{2}$.
\end{theorem}

\begin{proof}
Since $g=\frac{1}{2\pi}\int\gh(\omega)e^{i\omega\cdot}\dd\omega$ by inversion, and since
$\gh\in\Sw\subset\Lone$, Fubini gives
\begin{align}
\int_{\R}f(t)\overline{g(t)}\dd t
&=\int_{\R}f(t)\overline{\frac{1}{2\pi}\int_{\R}\gh(\omega)e^{i\omega t}\dd\omega}\dd t \notag\\
&=\frac{1}{2\pi}\int_{\R}\overline{\gh(\omega)}\Big[\int_{\R}f(t)e^{-i\omega t}\dd t\Big]\dd\omega
=\frac{1}{2\pi}\int_{\R}\fh(\omega)\overline{\gh(\omega)}\dd\omega .
\end{align}
Taking $g=f$ gives the norm identity. The extension: $\Sw$ is dense in $L^{2}$, and on $\Sw$ the map
$\F/\sqrt{2\pi}$ preserves the norm, so it extends uniquely by continuity to an isometry of $L^{2}$;
invertibility follows from \eqref{eq:F2}.
\end{proof}

\begin{corollary}\label{cor:no-energy}
The constant $1$ and the distribution $\delta$ are not elements of $L^{2}$, and neither can be
assigned a finite energy. For $f_{\ep}=e^{-\ep|t|}$,
\begin{equation}
\|f_{\ep}\|_{2}^{2}=\frac{1}{\ep}\longrightarrow\infty,
\qquad
\frac{1}{2\pi}\|G_{\ep}\|_{2}^{2}=\frac{1}{\ep},
\end{equation}
so the constant value $2\pi$ of the area of $G_{\ep}$ carries no information about
\eqref{eq:plancherel-pair} in the limit.
\end{corollary}

\begin{proof}
$\int_{\R}1\dd t=\infty$; for $f_{\ep}$ the two sides of
\eqref{eq:plancherel-pair} are computed by direct integration, using
$\int_{\R}(\ep^{2}+\omega^{2})^{-2}\dd\omega=\pi/(2\ep^{3})$.
\end{proof}

\begin{remark}[three layers, and their intersection]
On $L^{1}$ the transform is an absolutely convergent integral; on $L^{2}$ it is an isometry obtained
by density; on $\Sw'$ it is defined by duality and exists for objects---$\delta$, $1$, $t^{n}$,
$H$---that belong to neither $L^{1}$ nor $L^{2}$. The three definitions agree on their common domain,
which is $\Sw$ at least, and this agreement is the content of
Proposition~\ref{prop:dual-consistency}. The constructions do not replace one another; they are the
same operator seen in three topologies.
\end{remark}

\section{Summary of the chapter}

\begin{center}
\small
\begin{tabular}{@{}p{3.4cm}p{9.4cm}@{}}
\toprule
Step & Content \\
\midrule
Obstruction & $\F$ does not preserve $C_{c}^{\infty}$ (Proposition~\ref{prop:FnotD}) \\
Replacement & $\Sw$, closed under differentiation and multiplication (Definition~\ref{def:S}) \\
Closure & $\F:\Sw\to\Sw$ continuous (Theorem~\ref{prop:F-S}) \\
Inversion & proved by Gaussian regularization, with no appeal to $\delta$
(Theorem~\ref{thm:inversion}) \\
Definition & $\ip{\F T}{\varphi}=\ip{T}{\F\varphi}$ on $\Sw'$
(Definition~\ref{def:FTdist}) \\
Continuity & $T_{j}\to T$ implies $\F T_{j}\to\F T$ (Theorem~\ref{thm:F-cont}) \\
Harvest & $\F[\delta]=1$, $\F[1]=2\pi\delta$, $\F[t^{n}]=2\pi i^{n}\delta^{(n)}$,
$\F[H]=\pi\delta+\PV\frac{1}{i\omega}$ (Theorem~\ref{thm:main}) \\
Consistency & Gaussian and sinc families give the same limits \\
\bottomrule
\end{tabular}
\end{center}
